% ID: SEL-D0002-OB-E01
% ESTADO: cotejada. Fuente y transcripción contrastadas de forma independiente; salida revisada.
% Fuente: Selectividad/sel_2019_matematicas_aplicadas/Titular Septiembre/Exámenes Mat Apl CCSS II - A y B.pdf
% SHA-256: 31fb46947e141f5c8da8f7b381e7d27b67fe5a9d1597120454ce66fce0313f1a
% Página: 2; opción B; ejercicio 1.
% Transcripción de primera pasada basada en la imagen de la página original.
% Se conservan los datos y la construcción «el número ... debe ser» de la fuente.
% Sin solución ni modificación de datos o condiciones.

Una empresa comercializa dos tipos de concentrado de café, A y B, que se obtienen a partir de tres tipos de grano: de Colombia, de Etiopía y de Costa Rica. Para elaborar 1 kg de concentrado A se necesitan 4.5 kg de grano de Colombia y 3 kg de grano de Etiopía. Por otra parte, se requieren 7.5 kg de grano de Colombia y 1.5 kg de grano de Costa Rica para elaborar 1 kg de concentrado B. Actualmente la empresa dispone de un máximo de 67.5 kg de grano de Colombia, 30 kg de grano de Etiopía y 9 kg de grano de Costa Rica. Además, se exige que el número de kilogramos de concentrado A producidos debe ser mayor o igual que la mitad de los kilogramos de concentrado B.
\begin{enumerate}[label=\alph*),leftmargin=2.4em]
  \item \textbf{(1.75 puntos)} Represente la región factible que describe el problema anterior y determine sus vértices.
  \item \textbf{(0.25 puntos)} Indique de manera razonada si con las condiciones dadas sería posible producir 7 kg del concentrado A y 5 kg del concentrado B.
  \item \textbf{(0.5 puntos)} Sabiendo que el beneficio obtenido por la venta de cada kilogramo de concentrado del tipo A es 2 euros y de cada kilogramo del tipo B es 4 euros, \textquestiondown{}cuántos kilogramos del tipo A y cuántos del tipo B se habrán de producir para que el beneficio sea máximo? \textquestiondown{}Cuál es ese beneficio?
\end{enumerate}
